% LABEL -- os/samir/cmd/report.c lbl_run (FORM <file>, default .LBL, the
% 1034-byte dBASE III PLUS label form; TO FILE <file>, default .TXT; every
% record in master-index order; 54 for an invalid form).  Lane samirdml
% 14a8fc75.  Gates: test-samir-dml, test-samir-dml-emu.
\begin{XbaseEntry}{LABEL}{command}
Prints labels from a label form file.
\Syntax
\SyntaxLine{\kw{LABEL FORM} \xp{label file} \opt{\kw{TO FILE} \xp{file}}}
\Usage
\cmd{LABEL FORM} prints one label for every record of the database in USE,
in the order of the master index, laid out by the label form file named.
If no extension is given, \file{.LBL} is assumed.  The label form gives the
height and width of a label, the margin, the space between labels, the
number of labels across, and the expressions printed on each line of the
label.

\kw{TO FILE} writes the labels to a text file instead of the screen;
\file{.TXT} is assumed.

InitechBase prints label forms made with dBASE III PLUS; it does not create
them.  A form that is not a valid label form is reported as
\xmsg{54}{Label file invalid.}
\Examples
The label form \file{STAFFLBL.LBL} on the sample diskette prints three
interoffice mail labels across, each three lines high: the name, the
telephone extension and the department.  With \file{STAFFDPT.NDX} as the
master index the labels come in department order:
\DotExample{examples/c-label}
\SeeAlso{REPORT FORM}
\end{XbaseEntry}

